\chapter{Navier--Stokes: the terms equation (1) leaves out}

\textbf{Purpose:} Hold every claim of the first observation of the thought experiment beside what
stands behind it. A reader then sees at a glance which claims are proved, which were read, which
arrived by report, and which are still only theory. \textbf{Scope:} the thought experiment recorded
on 2026-10-01 as a run of questions and answers, and the research paper chapter \emph{The terms
equation (1) leaves out}, whose propositions this chapter names by number.

\section{How an entry is kept}

Every claim carries one status, and the status says what backs it.

\begin{center}
\begin{tabular}{@{}p{0.18\textwidth}>{\raggedright\arraybackslash}p{0.74\textwidth}@{}}
\toprule
Status & What it means \\
\midrule
proved & a proof is written out in the research paper and the number of its proposition is given \\
read & a source was opened and the claim is in it, in the words given \\
by report & the claim reached this work through another source, and that source was not opened \\
theory & stated, argued from scaling or by analogy, and not proved \\
not so & fails on its own terms; kept, with the reason \\
open & neither proved nor shown false here \\
\bottomrule
\end{tabular}
\end{center}

A claim changes status only when a proof or a reading changes it, and the change is recorded
beside it. Nothing is deleted. A claim that failed stays, with what failed it.

\section{The first observation}

The thought experiment opens by asking where the Navier--Stokes equations are bound: where a value
is assigned to the fluid instead of being asked of it. Its answer is that they are bound in nearly
every term. Density is assigned, viscosity is assigned, temperature is left out, and the fluid is
put in a box before anything is known about its container. Each of these is known from measurement
to be wrong for a real fluid, or is a choice made before the fluid was looked at.

The program it sets is to write every term in vector form and to hold nothing fixed that the fluid
has not been asked for. The fluid fills all of space until a container is measured, and the
container is read from the motion, through spherical harmonics of the wave patterns the motion
carries, and not assumed.

The same shape runs through the whole set of equations: put the fluid in a box, approximate what
it does there, and average the temperature over a volume. Each step drops information the motion
carries. This chapter enters the dropped terms and the box. Averaging over a volume is the second
observation and has the next chapter. Every other step of that shape goes in a later one.

The record goes on to entropy, observation and arithmetic. Those parts bear on no claim in this
observation and are not entered here.

\section{Entries}

\begin{longtable}{@{}>{\raggedright\arraybackslash}p{0.40\textwidth}>{\raggedright\arraybackslash}p{0.12\textwidth}>{\raggedright\arraybackslash}p{0.40\textwidth}@{}}
\toprule
Claim & Status & What stands behind it \\
\midrule
\endhead
\bottomrule
\endfoot
Equation (1) is bound: its constants are assigned to the fluid and not asked of it.
& read
& Fefferman's (1) carries one constant, $\nu$, and no density and no temperature. \\[0.6em]
Density varies with temperature, and a constant density is assigned instead of measured.
& proved
& Proposition 4: with $\rho'(T)\neq 0$, condition (2) forces conduction to cancel shear heating at
every point. Silent where $\rho'$ vanishes. \\[0.6em]
Viscosity is not a constant, and pulling $\mu$ out of the divergence drops a term.
& proved
& Propositions 1 and 3: the dropped term is $2\mu'(T)\,D\nabla T$, zero only where $\nabla T$ lies
in the null space of $D$. \\[0.6em]
The fluid defines its own state if it is let.
& proved
& Proposition 2: any viscous motion that decays at infinity makes its own temperature
gradient. \\[0.6em]
With vector calculus the dropped terms cannot be hidden.
& proved
& Proposition 1: the shear part is twelve terms per component and one product in vector
form. \\[0.6em]
Friction is not linear, anywhere.
& open
& The dependence measured in the research paper is on temperature, $-2.45$ percent per kelvin
for water at $20\,^\circ$C, evaluated from IAPWS R12-08~\cite{src:IAPWS-R12-08}. A dependence of water's
viscosity on $|D|$ was not measured. Read since: Stokes reports from the pipe experiments of
his time that the tangential force at a wall varies nearly as the square of the velocity, except at
small speeds~\cite{src:Stokes-1845}, and Navier takes the linear law as a principle. Neither
measures water's viscosity against $|D|$, and the status stays open. \\[0.6em]
Put a vector magnitude where a value is not known.
& read
& A viscosity $\mu(|D|) = \mu_0|D|^{p-2}$ is that move. For $p > 11/5$ in three dimensions a
solution obeying the energy equality is obtained; for $p = 2$ that is not known. Read in
Lienstromberg, Schiffer and Schubert~\cite{src:Lienstromberg-Schiffer-Schubert-2025}. \\[0.6em]
Never put the fluid in a box unless it is in one; a box destroys information.
& proved
& Proposition 6: the box keeps the transform of the field on the integer lattice and drops the
rest, and nonzero flows with zero divergence enter it as zero. It loses nothing when the fluid
already fits in one cell. \\[0.6em]
The box sets directions the fluid does not have.
& proved
& Proposition 7: 48 rotations and reflections about a point act on the box, against an infinite
group on all of space. \\[0.6em]
Spherical harmonics of the wave patterns recover the container.
& theory
& Stated in the thought experiment. Not proved, and no computation was run. \\[0.6em]
The forced construction breaks down under real conditions.
& proved, in part
& Proposition 5: the rate of shear heating has no bound at the singular time, and its total is
finite. That the breakdown would not happen with the coupling kept is open. \\[0.6em]
The forced construction failed to prove the theorem.
& not so
& The theorem is about equation (1), and its proof was not read. Nothing here bears on whether
the proof is correct. What the propositions bound is what the theorem is about. Read since, at the
level of its statements, its order of choices and its closing estimates; nothing read bears
against it~\cite{src:OpenAI-2026-Navier-Stokes}. \\[0.6em]
A collapsing vortex in water vaporizes near 0.7 nanometers of width, found by independent work.
& by report
& Reached this work only through the thought experiment. No source was located or opened, and
nothing rests on it. Since: the nearest source found puts the first failure in water at
cavitation near a millimeter of core radius and gives 0.3 nanometers as the molecular scale; its
0.7 is in micrometers and for air~\cite{src:Duraiswami-2026}. It bears on the claim and does not
support it. \\[0.6em]
An unbounded value is a sign that an assumption has left the fluid.
& theory
& Proposition 5 is the proved part: the unbounded factor is a factor of the dropped term. The
scaling entry below is the rest. Read since: of a light-emitting bubble, ``since the light emission
occurs as a (mathematical) singularity is forming, any transport process which a theorist chooses
to incorporate in this model can affect the result''~\cite{src:Barber-1997}; and of the forced
construction, a real fluid leaves the description of the equations long before the
singularity~\cite{src:Duraiswami-2026}. \\
\end{longtable}

\section{A scaling estimate, marked as theory}

Take a collapsing core of radius $r$ and speed $U$, with $|D| \sim U/r$. Suppose the shear heat in
the core, of order $\mu U^2/r^2$ per unit volume, is carried off by conduction across the core, of
order $k\,\Delta T/r^2$. Then the temperature excess of the core is
\[
\Delta T \sim \frac{\mu U^2}{k},
\]
independent of $r$. The ratio of the dropped term to the kept viscous term is of order
$|2D\nabla\mu| / |\mu\Delta u| \sim 2\,|(\ln\mu)'|\,\Delta T$, since both carry one more factor of
$1/r$ than $D$. The ratio is therefore
\[
2\,|(\ln\mu)'|\,\frac{\mu U^2}{k},
\]
and it passes one at
\[
U^* = \sqrt{\frac{k}{2\mu\,|(\ln\mu)'|}} .
\]
Any collapse in which $U$ grows without bound passes $U^*$. Below it the dropped term is a small
error. Above it the dropped term is as large as the term kept.

\textbf{Why this is theory and not proved.} It assumes one length scale for the core, a balance
between shear heating and conduction that holds at every moment, and derivatives estimated by size
alone. None of the three is shown.

For water at $20\,^\circ$C and atmospheric pressure, IAPWS R15-11~\cite{src:IAPWS-R15-11} gives $k = 0.5980$ watts per
meter per kelvin and IAPWS R12-08 gives $\mu = 1.0016$ mPa\,s with $(\ln\mu)' = -0.0245$ per
kelvin, and $U^* = 110$ meters per second. At $0\,^\circ$C it is 67 meters per second, and at
$40\,^\circ$C it is 160. The number is computed. The estimate it comes from is still theory.

\textbf{Read since.} Duraiswami~\cite{src:Duraiswami-2026} puts the axis of the forced
construction in water at cavitation when the swirl reaches 10 to 17 meters per second, at a core
radius near a millimeter, with the viscous heating over the rest of the collapse near 0.01 kelvin.
On those numbers water cavitates well below $U^*$, and the dropped term is not the first
assumption to fail in water: the pressure floor is. The estimate stays theory, and it now has a
competitor that fails first.

Gruntfest and Becker~\cite{src:Gruntfest-and-Becker-1964} find that shear at a fixed stress
between walls held at a fixed temperature, in a liquid whose viscosity falls exponentially with
temperature, has no steady state above a peak stress. Water at $20\,^\circ$C reaches that peak
when the walls move at 666 meters per second, $6.04\,U^*$. With the walls moving at a fixed speed
the steady layer exists at every speed, and its middle stands at $T_c$ with
$V^2 = 8\int_{T_0}^{T_c} k/\mu\,dT$, independent of the gap. On the IAPWS values the middle
reaches $100\,^\circ$C only near 956 meters per second, and at the 17 meters per second of
cavitation in the forced construction it stands 0.06 kelvin above the walls.

\section{What it wants}

Every term this chapter enters is collected, with the others, in the table of terms to put back
in the last chapter.

The estimate above wants a computation: the forced construction, or any collapsing solution of
(1), with $T$ carried beside it through the energy balance, and $2\mu'(T)\,D\nabla T$ evaluated
along the way and compared with $\mu\Delta u$. That replaces the scaling with a number taken on an
actual solution. It is not built.

The gap in Proposition 3 wants a proof either way: whether a flow can hold its own temperature
gradient inside the null space of its own rate of strain at every point and time.

The spherical harmonic entry wants a statement precise enough to prove: what feature of a
container, read from which harmonics of the motion on all of space, and with what loss.
